Angenommen werden eine Reaktionszeit $t_\mathrm{R} = \tReO$ und eine Bremsverzögerung auf trockener Strasse von $a = \aBrO$. Der Anhalteweg besteht aus dem Reaktionsweg (gleichförmig) und dem Bremsweg (gleichmässig verzögert bis zum Stillstand, mittlere Geschwindigkeit $v_1/2$ während der Bremszeit $v_1/a$):
\begin{align}
 s_\mathrm{R} &= \sRaF = \vAa\cdot\tRe = \sRa \\
 s_\mathrm{B} &= \frac{v_1}{2}\cdot\frac{v_1}{a} = \sBaF = \frac{\left(\vAa\right)^2}{2\cdot\aBr} = \sBa \\
 s_\mathrm{A} &= s_\mathrm{R}+s_\mathrm{B} = \sAn \approx \result{\sAnP{0}{2}}
\end{align}
Das zweite Auto fährt während der Reaktionszeit ungebremst
\begin{align}
 s_\mathrm{R,2} &= \sRbF = \vAb\cdot\tRe = \sRb
\end{align}
Das ist mehr als der ganze Anhalteweg des ersten Autos. Am Anhalteort des ersten Autos hat der Fahrer noch nicht einmal zu bremsen begonnen; das Auto fährt dort noch mit der vollen Geschwindigkeit von $\result{\vAbkP{0}{2}}$.

Das Ergebnis hängt von den angenommenen Werten ab, bleibt aber ähnlich: Auch mit einer etwas kürzeren Reaktionszeit wird das zweite Auto bis zu diesem Ort kaum langsamer. Sein eigener Anhalteweg beträgt $\sAbF = \sAbP{0}{2}$, also etwa das Doppelte.
